% META: {"id": "1SPE-VARALEA-CO-048", "chapitre": "1SPE-VARIABLES-ALEATOIRES", "type_objet": "corrige", "exercice_ref": "1SPE-VARALEA-EX-048", "capacites_codes": ["C5"], "sources_inspiration": [], "status": "generated"}
% BEGIN-VERIFY
% from sympy import Rational, sqrt
% p = Rational(3,5)
% E = 8000*p + (-10000)*(1-p)
% assert E == 800
% E2 = 8000**2*p + 10000**2*(1-p)
% V = E2 - E**2
% assert V == 77760000
% sigma = sqrt(V)
% assert round(float(sigma)) == 8818
% # 5 projets independants : Y = gain d'un projet, S = gain total
% E_Y = 1600*p + (-2000)*(1-p)
% assert E_Y == 160
% E_Y2 = 1600**2*p + 2000**2*(1-p)
% assert E_Y2 == 3136000
% V_Y = E_Y2 - E_Y**2
% assert V_Y == 3110400
% E5 = 5 * E_Y
% assert E5 == 800
% V_S = 5 * V_Y
% assert V_S == 15552000
% sigma_S = sqrt(V_S)
% assert abs(float(sigma_S) - 3943.602) < 1e-3
% assert round(float(sigma_S)) == 3944
% # la diversification divise l'ecart-type par sqrt(5)
% assert abs(float(sigma/sigma_S) - float(sqrt(5))) < 1e-9
% END-VERIFY
\begin{corrige}{1SPE-VARALEA-EX-048}

\textbf{1.} $G = 8000$ si succès ($p = 0{,}6$), $G = -10000$ si échec ($p = 0{,}4$).

\medskip
\textbf{2.} $E(G) = 8000 \times 0{,}6 + (-10000) \times 0{,}4 = 4800 - 4000 = 800$~euros. L'investissement est intéressant en moyenne.

\medskip
\textbf{3.} $E(G^2) = 64\,000\,000 \times 0{,}6 + 100\,000\,000 \times 0{,}4 = 38\,400\,000 + 40\,000\,000 = 78\,400\,000$.

$V(G) = 78\,400\,000 - 640\,000 = 77\,760\,000$. $\sigma(G) \approx 8\,818$~euros. Le risque est élevé.

\medskip
\textbf{4.} Avec $5$ projets à $2000$~euros : le gain $Y$ d'un projet vaut $1600$ (succès) ou $-2000$ (échec), donc $E(Y) = 1600 \times 0{,}6 + (-2000) \times 0{,}4 = 160$ et le gain total $S$ vérifie $E(S) = 5 \times 160 = 800$~euros : même espérance.

$E(Y^2) = 1600^2 \times 0{,}6 + 2000^2 \times 0{,}4 = 1\,536\,000 + 1\,600\,000 = 3\,136\,000$, donc $V(Y) = 3\,136\,000 - 160^2 = 3\,110\,400$.

Les cinq projets étant indépendants, $V(S) = 5 \times 3\,110\,400 = 15\,552\,000$ et $\sigma(S) = \sqrt{15\,552\,000} \approx 3\,944$~euros. L'écart type est exactement divisé par $\sqrt{5}$ par rapport au projet unique ($8\,818 / \sqrt{5} \approx 3\,944$) : le risque est réduit par la diversification.

\end{corrige}
